\documentclass[main.tex]{subfiles}

\title{Frullani's Theorem}
\author{Senan Sekhon}
\date{\DTMToday}
% Created September 15, 2021

\begin{document}

\maketitle

\section{Introduction}

\begin{theorem}[Frullani's Theorem]\label{frullanis_theorem}
    Suppose $f\in C[0,\infty)$, $\lim_{x\to\infty} f(x)=L$ and $a,b>0$. Then:
    \begin{equation*}
        \int_0^\infty \frac{f(bx)-f(ax)}{x}\dx=\qty(L-f(0))\log(\frac{b}{a})
    \end{equation*}
\end{theorem}
\begin{proof}
    % Adapted from https://math.stackexchange.com/a/61841
    If $a=b$, the result holds trivially as both sides are zero. Assume without loss of generality that $a<b$ (otherwise, switch $a$ and $b$). Suppose $u,v>0$ and $av>bu$. Then we have:
    \begin{align*}
        \int_u^v \frac{f(bx)-f(ax)}{x}\dx
        &= \int_u^v \frac{f(bx)}{x}\dx - \int_u^v \frac{f(ax)}{x}\dx \\
        &= \int_{bu}^{bv} \frac{f(x)}{x}\dx - \int_{au}^{av} \frac{f(x)}{x}\dx \\
        &= \int_{av}^{bv} \frac{f(x)}{x}\dx + \int_{bu}^{av} \frac{f(x)}{x}\dx - \int_{bu}^{av} \frac{f(x)}{x}\dx - \int_{au}^{bu} \frac{f(x)}{x}\dx \\
        &= \underbrace{\int_{av}^{bv} \frac{f(x)}{x}\dx}_{G(v)} - \underbrace{\int_{au}^{bu} \frac{f(x)}{x}\dx}_{G(u)}
    \end{align*}
    We now show that $\lim_{u\to0^+} G(u)=f(0)\log(\frac{b}{a})$. Suppose $\ep>0$. Then there exists $\delta>0$ such that $\abs{f(u)-f(0)}<\frac{\ep}{\log(b/a)}$ for all $0\le u<\frac{\delta}{b}$ (and so $au,bu<\delta$). This yields:
    \begin{align*}
        \abs{G(u)-f(0)\log(\frac{b}{a})}
        &= \abs{\int_{au}^{bu} \frac{f(x)}{x}\dx-f(0)\log(\frac{b}{a})}
        = \abs{\int_{au}^{bu} \frac{f(x)}{x}\dx-f(0)\int_{au}^{bu} \frac{1}{x}\dx} \\
        &= \abs{\int_{au}^{bu} \frac{f(x)-f(0)}{x}\dx}
        \le \int_{au}^{bu} \frac{\abs{f(x)-f(0)}}{x}\dx \\
        &\le \frac{\ep}{\log(b/a)}\int_{au}^{bu} \frac{1}{x}\dx
        = \frac{\ep}{\log(b/a)}\log(\frac{b}{a})
        =\ep
    \end{align*}
    Thus $\lim_{u\to0^+} G(u)=f(0)\log(\frac{b}{a})$.\\

    We now show that $\lim_{v\to\infty} G(v)=L\log(\frac{b}{a})$. Suppose $\ep>0$. Then there exists $H>0$ such that $\abs{f(x)-L}<\frac{\ep}{\log(b/a)}$ for all $x>H$. This yields:
    \begin{align*}
        \abs{G(v)-L\log(\frac{b}{a})}
        &= \abs{\int_{av}^{bv} \frac{f(x)}{x}\dx-L\log(\frac{b}{a})}
        = \abs{\int_{av}^{bv} \frac{f(x)}{x}\dx-L\int_{av}^{bv} \frac{1}{x}\dx} \\
        &= \abs{\int_{av}^{bv} \frac{f(x)-L}{x}\dx}
        \le \int_{av}^{bv} \frac{\abs{f(x)-L}}{x}\dx
        \le \frac{\ep}{\log(b/a)}\int_{av}^{bv} \frac{1}{x}\dx
        \le \frac{\ep}{\log(b/a)}\log(\frac{b}{a})
        =\ep
    \end{align*}
    Thus $\lim_{v\to\infty} G(v)=L\log(\frac{b}{a})$. Finally, we have:
    \begin{equation*}
        \int_0^\infty \frac{f(bx)-f(ax)}{x}\dx
        = \lim_{v\to\infty} G(v) - \lim_{u\to0^+} G(u)
        = (L-f(0))\log(\frac{b}{a}) \qedhere
    \end{equation*}
\end{proof}
\begin{proof}[A simpler proof, assuming $f$ is differentiable and $f'$ is integrable on $[0,\infty)$]
    \begin{equation*}
        \int_0^\infty \frac{f(bx)-f(ax)}{x}\dx
        = \int_0^\infty \int_a^b f'(xy)\dx\dy
        = \int_a^b \int_0^\infty f'(xy)\dy\dx
        = \int_a^b \frac{L-f(0)}{y}\dy
        = (L-f(0))\log(\frac{b}{a}) \qedhere
    \end{equation*}
\end{proof}

Frullani's theorem was first stated\footnote{\url{http://www.ams.org/journals/proc/1990-109-01/S0002-9939-1990-1007485-4/S0002-9939-1990-1007485-4.pdf}} in 1821 in a letter from Giuliano Frullani (1795--1834) to \href{https://mathshistory.st-andrews.ac.uk/Biographies/Plana/}{Giovanni Plana} (1781--1864), though a complete proof
was only given in 1823 by \href{https://mathshistory.st-andrews.ac.uk/Biographies/Cauchy/}{Augustin-Louis Cauchy} (1789--1857).

\begin{example}
    $f(x)=e^{-x}$
    \begin{equation*}
        \int_0^\infty \frac{e^{-bx}-e^{-ax}}{x}\dx = -\log(\frac{b}{a})
    \end{equation*}
\end{example}

\begin{example}
    $f(x)=\arctan(x)$
    \begin{equation*}
        \int_0^\infty \frac{\arctan(bx)-\arctan(ax)}{x}\dx = \frac{\pi}{2}\log(\frac{b}{a})
    \end{equation*}
\end{example}

\begin{example}
    $f(x)=\arcsec(x+1)$
    \begin{equation*}
        \int_0^\infty \frac{\arcsec(bx+1)-\arcsec(ax+1)}{x}\dx = \frac{\pi}{2}\log(\frac{b}{a})
    \end{equation*}
\end{example}

\begin{example}
    $f(x)=\tanh(x)$
    \begin{equation*}
        \int_0^\infty \frac{\tanh(bx)-\tanh(ax)}{x}\dx = \log(\frac{b}{a})
    \end{equation*}
    For certain values of $a$ and $b$, this can be rewritten in other forms, e.g.
    \begin{align*}
        \tanh(2x)-\tanh(x)
        &= \tanh(x)\sech(2x) \\
        \tanh(3x)-\tanh(x)
        &= \frac{2\tanh(x)\sech(x)^2}{1+3\tanh(x)^2} \\
        \tanh(3x)-\tanh(2x)
        &= \frac{\tanh(x)\sech(x)^4}{\qty\big(1+\tanh(x)^2)\qty\big(1+3\tanh(x)^2)}
    \end{align*}
\end{example}

\begin{example}
    $f(x)=\sech(x)$
    \begin{equation*}
        \int_0^\infty \frac{\sech(bx)-\sech(ax)}{x}\dx = -\log(\frac{b}{a})
    \end{equation*}
    For certain values of $a$ and $b$, this can be rewritten in other forms, e.g.
    \begin{align*}
        \sech(x)-\sech(2x)
        &= \frac{\sech(x)(1-\sech(x))(2+\sech(x))}{2-\sech(x)^2} \\
        \sech(x)-\sech(3x)
        &= \frac{4\sech(x)\tanh(x)^2}{4-3\sech(x)^2}
    \end{align*}
\end{example}

\begin{example}
    $f(x)=\arctan(\sinh(x))$
    \begin{equation*}
        \int_0^\infty \frac{\arctan(\sinh(bx))-\arctan(\sinh(ax))}{x}\dx = \frac{\pi}{2}\log(\frac{b}{a})
    \end{equation*}
    \begin{align*}
        \arctan(\sinh(2x))-\arctan(\sinh(x))
        &=\arctan(\frac{\sinh(2x)-\sinh(x)}{1+\sinh(x)\sinh(2x)}) \\
        &=\arctan(\frac{2\sinh(x)\cosh(x)-\sinh(x)}{1+\sinh(x)\qty(2\sinh(x)\cosh(x))}) \\
        &=\arctan(\frac{\sinh(x)(2\cosh(x)-1)}{1+2\sinh(x)^2\cosh(x)})
    \end{align*}
\end{example}

The last part of the proof of \hyperref[frullanis_theorem]{Frullani's theorem} can be adapted to obtain a more general result:

\begin{theorem}
    % https://www.pnas.org/doi/pdf/10.1073/pnas.35.10.612
    Suppose $f\in C[0,\infty)$, $\lim_{x\to\infty} \frac{1}{x}\int_0^x f(y)\dy=L$ and $a,b>0$. Then:
    \begin{equation*}
        \int_0^\infty \frac{f(bx)-f(ax)}{x}\dx=\qty(L-f(0))\log(\frac{b}{a})
    \end{equation*}
\end{theorem}

\begin{example}
    $f(x)=\cos(x)$
    \begin{equation*}
        \int_0^\infty \frac{\cos(bx)-\cos(ax)}{x}\dx = -\log(\frac{b}{a})
    \end{equation*}
    In this case, $\lim_{x\to\infty} \frac{1}{x}\int_0^x \cos(y)\dy=\lim_{x\to\infty} \frac{\sin(x)}{x}=0$ while $\lim_{x\to\infty} \cos(x)$ does not exist.\\
    
    Another way to evaluate this integral is to use Frullani's theorem with $f(x)=e^{-px}\cos(x)$ for $p>0$, then take the limit as $p\to0^+$. However, this requires care as the integral is not absolutely convergent for $p=0$.
\end{example}
\begin{example}
    Suppose $p,q\in\R$, $p>\abs{q}$. Then we have:
    \begin{equation*}
        \int_0^\infty \frac{\sin(px)\sin(qx)}{x}\dx
        = \frac{1}{2}\int_0^\infty \frac{\cos((p-q)x)-\cos((p+q)x)}{x}\dx
        = \frac{1}{2}\log(\frac{p+q}{p-q})
        = \arctanh(\frac{q}{p})
    \end{equation*}
\end{example}



\begin{equation*}
    \int_0^\infty \frac{\sech(x)(1-\sech(x))(2+\sech(x))}{x\qty(2-\sech(x)^2)}\dx=\log(2)
\end{equation*}
\vspace{10pt}
\begin{equation*}
    \int_0^\infty \frac{1}{x}\arctan(\frac{\sinh(x)\qty(2\cosh(x)-1)}{1+2\sinh(x)^2\cosh(x)})\dx=\frac{\pi}{2}\log(2)
\end{equation*}
\vspace{10pt}
\begin{equation*}
    \int_0^\infty \frac{1}{x}\arccos(\frac{\cosh(4x)-\cosh(2x)+2}{2\cosh(x)\cosh(3x)})\dx=\frac{\pi}{2}\log(3)
\end{equation*}
\vspace{10pt}
\begin{equation*}
    \int_0^1 \frac{x}{(1+3x^2)\log(\dfrac{1+x}{1-x})}\dx=\frac{\log(3)}{4}
\end{equation*}

More examples are available at \url{https://scientia.mat.utfsm.cl/archivos/vol19/vol19art14.pdf}.

\begin{theorem}
    % https://www.pnas.org/doi/pdf/10.1073/pnas.35.10.612
    Suppose $f\in C[0,\infty)$, $\lim_{x\to\infty} \frac{1}{x}\int_0^x f(y)\dy=L$ and $a,b>0$. Then:
    \begin{equation*}
        \int_0^\infty \frac{f(bx)-f(ax)}{x}\dx=\qty(L-f(0))\log(\frac{b}{a})
    \end{equation*}
    Sub $u=e^x$ and $g(u)=f(\log(u))$, $\frac{1}{x}\int_1^{e^x} g(v)/v\dv=L$
    \begin{equation*}
        \int_1^\infty \frac{g(u^b)-g(u^a)}{u\log(u)}\du=\qty(L-f(0))\log(\frac{b}{a})
    \end{equation*}
\end{theorem}


\end{document}